\chapter{The Two Welfare Theorems}\label{ch:b2:welfare-theorems}

In 1855--1858, American red wheat sold in London for 22.0 per cent more on average than in New York.
% ledger: b2:persson2002-t3
In 1892--1900 the gap was 13.9 per cent. Part of a gap between two markets is the cost of moving the
good between them: freight, insurance, port charges. The chart splits each period's average gap into
those costs and a residual, the part that the cost of shipping does not explain.

\begin{openingchart}{The London--New York price gap for American wheat, in per cent of the New York
  price, in four periods, split into freight, insurance and port charges (an assumed 3.75 per cent
  before 1865 and 2.75 per cent after) and the residual.}{Persson (2002), Table~3}
\begin{tikzpicture}
\begin{axis}[omrecord, ybar stacked, bar width=26pt, xmin=0.4, xmax=4.6, ymin=0, ymax=30, xtick=data,
    xticklabels from table={figdata/bachelor-2/07-welfare-theorems/out/gap.csv}{period},
    ylabel={Per cent of the New York price}, grid=none, ymajorgrids, width=0.7\linewidth,
    height=0.48\linewidth, legend columns=3,
    legend style={at={(0.5,-0.16)}, anchor=north, font=\footnotesize}]
  \addplot[fill=omDef!70, draw=omDef] table[x=pos, y=transport]
    {figdata/bachelor-2/07-welfare-theorems/out/gap.csv};
  \addplot[fill=black!25, draw=black!50] table[x=pos, y=charges]
    {figdata/bachelor-2/07-welfare-theorems/out/gap.csv};
  \addplot[fill=omThm!60, draw=omThm] table[x=pos, y=residual]
    {figdata/bachelor-2/07-welfare-theorems/out/gap.csv};
  \legend{Freight, Insurance and port, Residual}
\end{axis}
\end{tikzpicture}
\end{openingchart}

% figdata: gap_summary (fall 8.1, transport_fall 2.3, share_transport 28, residual_first 8.4, residual_last 3.5)
From the first period to the last, the gap fell by 8.1 points. Freight fell by only 2.3 points, 28 per
cent of the fall; the residual fell from 8.4 to 3.5 points. A model of competitive markets predicts that
when a good is shipped from one market to another, the gap between their prices equals the cost of
shipping it: any larger gap is a profit that traders would compete away. Measured against that
prediction, the markets of the 1850s left a large unexplained margin, and the markets of the 1890s a
much smaller one. Persson concludes that better market and marketing efficiency did more than cheaper
freight to bring the two prices together.

That prediction is one consequence of a broader result. \Cref{ch:b2:exchange-economies} found competitive
equilibria in economies where people only trade what they hold. This chapter adds production and proves
the two theorems that relate \omterm{def:b2:exchange-economies:competitive-equilibrium}{competitive equilibrium} to Pareto efficiency: every \omterm{def:b2:exchange-economies:competitive-equilibrium}{competitive equilibrium}
is efficient, and, under more assumptions, every efficient \omterm{def:b2:exchange-economies:allocation}{allocation} is a \omterm{def:b2:exchange-economies:competitive-equilibrium}{competitive equilibrium} once
wealth has been redistributed. Each assumption the proofs use will be named, because each one fails
somewhere, and those failures fill much of the rest of this book.

\begin{recall}
From \cref{ch:b2:exchange-economies}: exchange economies, \omterm{def:b2:exchange-economies:competitive-equilibrium}{competitive equilibrium}, \omterm{def:b2:exchange-economies:walras-law}{Walras' law}, Pareto
efficiency and the \omterm{def:b2:exchange-economies:contract-curve}{contract curve}. From the Year 1 volume: the firm's profit maximisation, the
production possibility frontier, and consumer and producer surplus. From the school volumes: total
surplus as a measure of efficiency in one market.
\end{recall}

\begin{notation}[Symbols of this chapter]\label{not:b2:welfare-theorems:symbols}
As in \cref{ch:b2:exchange-economies}, consumers are $i$, goods $\ell$, \omterm{def:b2:exchange-economies:endowment}{endowments} $\omega_i$, bundles
$x_i$, and $p = (p_1, \dots, p_n)$ is a price vector; $p \cdot x = \sum_\ell p_\ell x_\ell$. Firms are $j$,
with production plans $y_j$ (inputs negative, outputs positive) chosen from technology sets $Y_j$;
$\theta_{ij}$ is consumer $i$'s share of firm $j$'s profit. In the one-consumer economy, $T$ is the time
available, $L$ hours worked, $c$ consumption and $w$ the real wage. $\tau_i$ is a \omterm{def:b2:welfare-theorems:lump-sum}{lump-sum transfer} to
consumer $i$, and $t$ a transport cost per unit.
\end{notation}

\section{Production added}\label{sec:b2:welfare-theorems:production}

\begin{definition}[Private-ownership economy]\label{def:b2:welfare-theorems:production-economy}
A \emph{private-ownership economy}\index{private-ownership economy} has consumers with preferences and
\omterm{def:b2:exchange-economies:endowment}{endowments}, firms $j$ with technology sets $Y_j \subset \mathbb{R}^n$, and ownership shares $\theta_{ij}
\ge 0$ with $\sum_i \theta_{ij} = 1$ for each firm. At prices $p$, each firm chooses $y_j \in Y_j$ to
maximise its profit $p \cdot y_j$, and each consumer chooses the best bundle with $p \cdot x_i \le p \cdot
\omega_i + \sum_j \theta_{ij}\,p \cdot y_j$. A \omterm{def:b2:exchange-economies:competitive-equilibrium}{competitive equilibrium} is a price vector and plans such
that every market clears: $\sum_i x_i = \sum_i \omega_i + \sum_j y_j$.
\end{definition}

The simplest case has one consumer, who owns one firm and works for it: the \emph{Robinson Crusoe
economy}. Crusoe has $T = 16$ hours a day, likes consumption and leisure, $u = \ln c + \ln(T - L)$, and
produces $c = 4\sqrt L$ with $L$ hours of work. A planner who chooses $L$ directly maximises $\ln(4\sqrt L)
+ \ln(16 - L)$; the first-order condition $1/(2L) = 1/(16 - L)$ gives $L = 16/3 = 5.333$ and $c = 4\sqrt{5.333}
= 9.238$.

% code: exercises/solutions.py crusoe
Now split Crusoe in two: a firm that hires labour at the real wage $w$ and sells output at a price of
one, and a household that sells labour and receives the firm's profit. The firm maximises $4\sqrt L - wL$,
so $w = 2/\sqrt L$. At $w = 0.866$ it hires $5.333$ hours, produces $9.238$ and makes a profit of $9.238 -
0.866 \times 5.333 = 4.619$. The household's income, valued at full time, is $16w + 4.619 = 18.475$; with
log utility it spends half on consumption and half on leisure, so it consumes $9.238$ and takes $18.475/
(2 \times 0.866) = 10.667$ hours of leisure: it works $5.333$ hours. The market outcome is the planner's.

\begin{omfigure}
\begin{tikzpicture}[econ]
\begin{axis}[omrecord, grid=none, xmin=0, xmax=17, ymin=0, ymax=17, xtick={0,4,8,12,16},
    ytick={0,4,8,12,16}, xlabel={Leisure $T - L$ (hours)}, ylabel={Consumption $c$},
    width=0.58\linewidth, height=0.55\linewidth]
  \addplot[omDef, very thick, domain=0:16, samples=120] {4*sqrt(16 - x)};
  \addplot[black, thick, domain=1:16, samples=2] {9.2376 - 0.8660*(x - 10.6667)};
  \addplot[omThm, thick, domain=5.9:16.5, samples=80] {98.535/x};
  \node[eqpoint] at (axis cs:10.667,9.238) {};
  \node[anchor=south west, font=\footnotesize] at (axis cs:10.8,9.4) {equilibrium};
  \node[eqpoint] at (axis cs:16,4.619) {};
  \node[font=\footnotesize] (pr) at (axis cs:11.5,2.0) {profit $4.619$};
  \draw[->, black!70] (pr.east) -- (axis cs:15.8,4.4);
  \node[anchor=west, font=\footnotesize, omDef] at (axis cs:0.6,6.0) {frontier $c = 4\sqrt{L}$};
  \node[anchor=south, font=\footnotesize] at (axis cs:4.0,12.0) {slope $-w$};
\end{axis}
\end{tikzpicture}

{\small The Robinson Crusoe economy. The frontier gives the consumption that each amount of leisure
leaves. At the real wage $w = 0.866$ the firm picks the point of the frontier where profit is highest;
the household's budget line has the same slope and meets the axis of full leisure at the profit it
receives. The household's indifference curve touches the same line at the same point.}
\end{omfigure}

The decentralised outcome works because the same price line is tangent to the frontier, where it
maximises the firm's profit, and to the indifference curve, where it maximises the household's utility.
Both facts will reappear in the proofs below, and the second will fail when the frontier is not convex.

\section{The first welfare theorem}\label{sec:b2:welfare-theorems:first}

\begin{definition}[Local non-satiation]\label{def:b2:welfare-theorems:lns}
Preferences are \emph{locally non-satiated}\index{local non-satiation} if for every bundle $x_i$ and
every distance $\varepsilon > 0$ there is a bundle within $\varepsilon$ of $x_i$ that the consumer
strictly prefers.
\end{definition}

Monotone preferences are \omterm{def:b2:welfare-theorems:lns}{locally non-satiated}: a little more of everything is better. The assumption
rules out a bliss point and \emph{thick} indifference curves, regions in which the consumer does not care
what she gets.

\begin{definition}[First fundamental theorem of welfare economics]\label{def:b2:welfare-theorems:first-theorem}
The \emph{first fundamental theorem of welfare economics}\index{first fundamental theorem of welfare
economics} is the result that every \omterm{def:b2:exchange-economies:competitive-equilibrium}{competitive equilibrium} \omterm{def:b2:exchange-economies:allocation}{allocation} is \omterm{def:b2:exchange-economies:pareto-efficient}{Pareto-efficient} when
preferences are \omterm{def:b2:welfare-theorems:lns}{locally non-satiated} (\cref{thm:b2:welfare-theorems:first}).
\end{definition}

\begin{theorem}[First welfare theorem]\label{thm:b2:welfare-theorems:first}
In an \omterm{def:b2:exchange-economies:exchange-economy}{exchange economy} with \omterm{def:b2:welfare-theorems:lns}{locally non-satiated} preferences, if $(p, x^*)$ is a \omterm{def:b2:exchange-economies:competitive-equilibrium}{competitive equilibrium}
with $p \gg 0$, then $x^*$ is \omterm{def:b2:exchange-economies:pareto-efficient}{Pareto-efficient}. The same holds in a \omterm{def:b2:welfare-theorems:production-economy}{private-ownership economy} with
production.
\end{theorem}

\begin{proof}
\omterm{def:b2:exchange-economies:exchange-economy}{Exchange economy}. Suppose a \omterm{def:b2:exchange-economies:feasible}{feasible allocation} $x'$ makes everyone at least as well off as $x^*$ and
consumer $k$ strictly better off. Since $x^*_k$ was the best bundle $k$ could afford, $x'_k$ was not
affordable: $p \cdot x'_k > p \cdot \omega_k$. For every other consumer, $p \cdot x'_i \ge p \cdot
\omega_i$: if instead $p \cdot x'_i < p \cdot \omega_i$, by \omterm{def:b2:welfare-theorems:lns}{local non-satiation} some bundle close
enough to $x'_i$ to remain affordable would be strictly preferred to $x'_i$, hence to $x^*_i$, which
contradicts the choice of $x^*_i$. Summing over consumers, $p \cdot \sum_i x'_i > p \cdot \sum_i
\omega_i$. But $x'$ is feasible, $\sum_i x'_i \le \sum_i \omega_i$, and with $p \gg 0$ this gives $p \cdot
\sum_i x'_i \le p \cdot \sum_i \omega_i$: a contradiction. The case with production is
\cref{exo:b2:welfare-theorems:10}.
\end{proof}

The proof is the figure below in symbols. At equilibrium prices, everything a consumer prefers costs
more than her wealth; adding up, any \omterm{def:b2:exchange-economies:allocation}{allocation} everyone prefers costs more than the economy has.

\begin{omfigure}
\begin{tikzpicture}[econ]
\begin{axis}[omrecord, grid=none, axis lines=box, xmin=0, xmax=10, ymin=0, ymax=10, xtick={0,2,4,6,8,10},
    ytick={0,2,4,6,8,10}, xlabel={Good 1 held by $A$}, ylabel={Good 2 held by $A$},
    width=0.58\linewidth, height=0.58\linewidth, clip=true]
  \addplot[omDef, thick, domain=3.05:10, samples=60] {30.4511/x};
  \addplot[omThm, thick, domain=0:9.2, samples=60] {10 - 98.9315/(10 - x)^2};
  \addplot[black, thick, domain=2.2:9.47, samples=2] {2 - 1.357143*(x - 8)};
  \node[eqpoint] at (axis cs:8,2) {};
  \node[anchor=north east, font=\footnotesize] at (axis cs:8,2) {$\omega$};
  \node[eqpoint] at (axis cs:4.737,6.429) {};
  \node[anchor=north east, font=\footnotesize] at (axis cs:4.6,6.3) {$x^*$};
  \node[font=\footnotesize, omDef, align=center] at (axis cs:7.6,7.6) {$A$ prefers:\\costs more than $p \cdot \omega_A$};
  \node[font=\footnotesize, omThm, align=center] at (axis cs:2.4,2.6) {$B$ prefers:\\costs $B$ more\\than $p \cdot \omega_B$};
\end{axis}
\end{tikzpicture}

{\small The first welfare theorem in the box of \cref{ch:b2:exchange-economies}. The budget line through
$\omega$ at the equilibrium price separates the bundles $A$ prefers to $x^*$ (above her indifference
curve) from those $B$ prefers (below his, seen from $A$'s corner). No point lies in both regions, so no
\omterm{def:b2:exchange-economies:allocation}{allocation} makes both better off.}
\end{omfigure}

The theorem extends to markets in different places. Treat wheat in New York and wheat in London as two
goods, and shipping as a firm that turns one unit of the first into one unit of the second at a cost $t$
per unit.

\begin{proposition}[Law of one price with transport costs]\label{prop:b2:welfare-theorems:one-price}
In a \omterm{def:b2:exchange-economies:competitive-equilibrium}{competitive equilibrium} with shipping firms that face a constant cost $t$ per unit, the prices of a
good in two markets satisfy $p_B - p_A \le t$ and $p_A - p_B \le t$, and $p_B - p_A = t$ whenever a
positive quantity is shipped from $A$ to $B$.
\end{proposition}

\begin{proof}
A shipping firm that buys $s \ge 0$ units in $A$ and sells them in $B$ earns $s(p_B - p_A - t)$. If $p_B -
p_A > t$, profit has no maximum: the firm would ship without limit, and the market for the good in $A$
could not clear. So $p_B - p_A \le t$, and likewise in the other direction. If $s > 0$ in equilibrium,
profit maximisation requires $p_B - p_A - t \ge 0$, so the inequality holds with equality.
\end{proof}

\begin{record}[Shipping and the law of one price]
Persson (2002) assembled weekly wheat prices in New York and in British markets and the freight between
% ledger: b2:persson2002-t1
them. Freight fell as a share of the New York wheat price, from 11.0 per cent in 1852--1860 to 7.5 per cent
in 1891--1900, a fall of almost a third (the author's own 23 per cent runs to his last, shorter period, 1896--1900,
at 8.5 per cent); but the gap between the London and New York prices of comparable
% ledger: b2:persson2002-t3, b2:persson2002-summary
wheat always exceeded freight and the assumed insurance and port charges. The residual was 8.4 per cent
of the New York price in 1855--1858 and 3.5 per cent in 1892--1900. The paper finds that comparing
different qualities of wheat in the two markets can create spurious convergence or divergence, and
concludes that improved market and marketing efficiency mattered more than freight.
\end{record}

The proposition gives the residual a meaning: in the model it is zero, so a positive residual is either a
cost the measure leaves out (the merchant's time, the risk of a price change during the voyage, the
interest on the cargo) or a profit that competition had not yet removed. Its fall says that one or the
other shrank. Which it was the data cannot say on their own.

\section{The second welfare theorem}\label{sec:b2:welfare-theorems:second}

The first theorem says markets reach an efficient \omterm{def:b2:exchange-economies:allocation}{allocation}, but which one depends on the \omterm{def:b2:exchange-economies:endowment}{endowments}.
The second asks the converse: can any efficient \omterm{def:b2:exchange-economies:allocation}{allocation} be reached by markets, if wealth is first
redistributed?

\begin{definition}[Lump-sum transfer]\label{def:b2:welfare-theorems:lump-sum}
A \emph{lump-sum transfer}\index{lump-sum transfer} is a transfer of wealth $\tau_i$ to consumer $i$,
with $\sum_i \tau_i = 0$, whose amount does not depend on any choice the consumer makes: not on what
she buys, sells, earns or works.
\end{definition}

\begin{definition}[Second fundamental theorem of welfare economics]\label{def:b2:welfare-theorems:second-theorem}
The \emph{second fundamental theorem of welfare economics}\index{second fundamental theorem of welfare
economics} is the result that, with convex preferences and technologies, every \omterm{def:b2:exchange-economies:pareto-efficient}{Pareto-efficient
allocation} can be reached as a \omterm{def:b2:exchange-economies:competitive-equilibrium}{competitive equilibrium} after suitable \omterm{def:b2:welfare-theorems:lump-sum}{lump-sum transfers}
(\cref{thm:b2:welfare-theorems:second}).
\end{definition}

\begin{theorem}[Second welfare theorem]\label{thm:b2:welfare-theorems:second}
In a \omterm{def:b2:welfare-theorems:production-economy}{private-ownership economy}, let preferences be continuous, convex and strictly monotone, and every
technology set $Y_j$ convex. Let $x^*$ be a \omterm{def:b2:exchange-economies:pareto-efficient}{Pareto-efficient allocation}, with production plans $y^*_j$,
in which every consumer consumes a positive amount of every good. Then there are prices $p \gg 0$ and
\omterm{def:b2:welfare-theorems:lump-sum}{lump-sum transfers} $\tau_i$ such that, with wealth $p \cdot \omega_i + \sum_j\theta_{ij}\,p \cdot y^*_j +
\tau_i$, each consumer's best bundle is $x^*_i$ and each firm's most profitable plan is $y^*_j$.
\end{theorem}

\admitted

The missing step is geometric. The set of aggregate bundles that could make everyone at least as well off
as at $x^*$, and someone better off, is convex, because preferences are; the set of bundles the economy
can produce is convex, because technologies are; efficiency says the two sets do not overlap. Two convex
sets that do not overlap can be separated by a hyperplane, and the hyperplane's normal vector is the
price vector $p$. The transfers then give each consumer exactly the wealth $p \cdot x^*_i$.

With Cobb--Douglas consumers the transfer can be computed. Take the running economy of
\cref{ch:b2:exchange-economies}, in which the \omterm{def:b2:exchange-economies:competitive-equilibrium}{competitive equilibrium} gives $A$ 4.737 units of good 1, and
suppose a planner wants $A$ to have 4. The listing finds the point of the \omterm{def:b2:exchange-economies:contract-curve}{contract curve}, the price that
supports it ($A$'s marginal rate of substitution there), and the transfer of good 2 that makes $A$'s wealth
at that price equal to the cost of her target bundle.

\omcode{code/bachelor-2/07-welfare-theorems/listings/second_theorem.py}{1}{29}{The transfer that turns a
chosen efficient \omterm{def:b2:exchange-economies:allocation}{allocation} into a \omterm{def:b2:exchange-economies:competitive-equilibrium}{competitive equilibrium}.}\label{lst:b2:welfare-theorems:second}

\omoutput{code/bachelor-2/07-welfare-theorems/outputs/second_theorem.txt}{A transfer of two units of good
2 from $A$ to $B$ moves the equilibrium to the target.}\label{out:b2:welfare-theorems:second}

% code: listings/second_theorem.py
At the price $1.4286 = 10/7$, $A$'s \omterm{def:b2:exchange-economies:endowment}{endowment} is worth $94/7$ and her target bundle costs $80/7$; taking
two units of good 2 from her closes the gap. Recomputing the equilibrium from the new \omterm{def:b2:exchange-economies:endowment}{endowments} returns
the same price and the target bundle: the transfer, not a change in how markets work, chose the
\omterm{def:b2:exchange-economies:allocation}{allocation}.

Convexity is where the theorem can fail. Suppose Crusoe's technology has increasing returns, $c =
L^2/10$, with $T = 10$ hours and $u = \ln c + \ln(T - L)$. The planner's condition $2/L = 1/(10 - L)$ gives
% code: exercises/solutions.py nonconvex
$L = 6.667$ and $c = 4.444$, an efficient plan. The price that would support it is the slope of the
frontier there, $w = 2L/10 = 1.333$. But at that wage the firm's profit from the efficient plan is $4.444 -
1.333 \times 6.667 = -4.444$, and larger plans earn more: with increasing returns, profit at a fixed wage
grows without bound. No wage makes the efficient plan the firm's choice.

\begin{omfigure}
\begin{tikzpicture}[econ]
\begin{axis}[omrecord, grid=none, xmin=0, xmax=10.5, ymin=0, ymax=10.5, xtick={0,2,4,6,8,10},
    ytick={0,2,4,6,8,10}, xlabel={Leisure $T - L$ (hours)}, ylabel={Consumption $c$},
    width=0.55\linewidth, height=0.55\linewidth]
  \addplot[omDef, very thick, domain=0:10, samples=100] {(10 - x)^2/10};
  \addplot[omThm, thick, domain=1.45:10.5, samples=80] {14.815/x};
  \addplot[black, thick, domain=0:6.67, samples=2] {4.4444 - 1.3333*(x - 3.3333)};
  \node[eqpoint] at (axis cs:3.333,4.444) {};
  \node[anchor=south west, font=\footnotesize] at (axis cs:3.5,4.6) {efficient plan};
  \node[eqpoint] at (axis cs:0,10) {};
  \node[anchor=west, font=\footnotesize, align=left] (mp) at (axis cs:4.5,9.2) {more profitable\\at the same wage};
  \draw[->, black!70] (mp.west) -- (axis cs:0.25,9.95);
  \node[anchor=north east, font=\footnotesize, omDef] at (axis cs:9.6,1.1) {frontier};
\end{axis}
\end{tikzpicture}

{\small A non-convex technology, $c = L^2/10$. The efficient plan is where the indifference curve touches
the frontier. The line through it with the frontier's slope is the only candidate price line, but the
frontier lies above it elsewhere: at that wage the firm would rather work all ten hours, and the plan
cannot be a \omterm{def:b2:exchange-economies:competitive-equilibrium}{competitive equilibrium}.}
\end{omfigure}

\section{Which assumption carries the result}\label{sec:b2:welfare-theorems:assumptions}

The first theorem used three facts: each consumer's choice at given prices, each firm's profit
maximisation at given prices, and that nothing affects a consumer except her own bundle. Hidden in the
statement is a fourth: there is a market, and a price, for everything consumers and firms care about.

\begin{definition}[Complete markets]\label{def:b2:welfare-theorems:complete-markets}
Markets are \emph{complete}\index{complete markets} if every good that enters a utility function or a
technology, at every date, in every place and in every \omterm{def:b2:insurance-and-risk-sharing:state}{state of the world}, can be traded at a price that
every agent takes as given.
\end{definition}

Each assumption fails somewhere, and each failure is a chapter of this book or the next.
\begin{itemize}
  \item \emph{No external effects.} When one agent's choice enters another's utility or technology
        without a price, as with pollution, the equilibrium is generally inefficient
        (\cref{ch:b2:externalities-and-public-goods-theory}).
  \item \emph{\omterm{def:b2:welfare-theorems:complete-markets}{Complete markets} and symmetric information.} When the parties to a trade know different
        things, some markets shrink or vanish (\cref{ch:b2:adverse-selection}), and contracts must trade off
        insurance against incentives (\cref{ch:b2:moral-hazard}). Missing markets for risk were the theme of
        \cref{ch:b2:insurance-and-risk-sharing}.
  \item \emph{Price taking.} A firm or consumer large enough to move prices does not trade where marginal
        values are equal; the strategic analysis begins in \cref{ch:b2:static-games}.
  \item \emph{Convexity.} The first theorem does not need it; the second does. Increasing returns, which
        are central to the economics of ideas, return in \cref{ch:b2:endogenous-growth}.
  \item \emph{\omterm{def:b2:welfare-theorems:lns}{Local non-satiation}.} Without it, an equilibrium can leave gains unclaimed
        (\cref{exo:b2:welfare-theorems:11}).
\end{itemize}
A failure of an assumption shows that the theorem does not apply. It does not show that some policy would
do better: correcting a market outcome needs information about preferences and technologies that the
market used without anyone collecting it, and the agencies that would correct it face their own
incentives. Both sides of that comparison are studied in the chapters just named and in the Year 3
volume.

\begin{model}{The competitive economy with production}
\assumes Price taking by every agent; \omterm{def:b2:welfare-theorems:lns}{locally non-satiated} preferences; no external effects; \omterm{def:b2:welfare-theorems:complete-markets}{complete
markets}; for the second theorem, convex preferences and technologies.
\predicts \omterm{def:b2:exchange-economies:pareto-efficient}{Pareto-efficient allocations}; profit-maximising production at prices equal to marginal rates
of transformation; price gaps between connected markets equal to the cost of moving goods between them.
\tested Commodity markets linked by trade: price gaps narrowed as freight and trading costs fell, and
the part of the gap that freight does not explain shrank in the second half of the nineteenth century
(Persson 2002); laboratory markets converge to competitive prices (\cref{ch:b2:exchange-economies}).
\breaks External effects, market power, missing markets and asymmetric information, and non-convex
technologies; and the theorems say nothing about which efficient \omterm{def:b2:exchange-economies:allocation}{allocation} is reached.
\end{model}

\section{The theorems are silent on distribution}\label{sec:b2:welfare-theorems:distribution}

The Pareto set is large: in the \omterm{def:b2:exchange-economies:edgeworth-box}{Edgeworth box} it runs from one corner to the other, from \omterm{def:b2:exchange-economies:allocation}{allocations} that
give $A$ almost nothing to \omterm{def:b2:exchange-economies:allocation}{allocations} that give $B$ almost nothing. The first theorem says that the
market picks one of them, the one that the \omterm{def:b2:exchange-economies:endowment}{endowments} lead to. It says nothing about whether that
\omterm{def:b2:exchange-economies:allocation}{allocation} is fair, or better than another efficient \omterm{def:b2:exchange-economies:allocation}{allocation}. Ranking efficient \omterm{def:b2:exchange-economies:allocation}{allocations} needs a
judgement about how to weigh one person's gain against another's loss, which economics can describe
(\cref{ch:b2:social-choice}) but not supply.

The second theorem seems to separate the two questions: choose the distribution by \omterm{def:b2:welfare-theorems:lump-sum}{lump-sum transfers},
then let markets deliver efficiency. The difficulty is the word \emph{lump-sum}. A transfer that does not
depend on anything a person does must be based on characteristics she cannot change, and those are the
characteristics a government cannot observe: ability, health or tastes, as opposed to income, spending or
hours, which respond to the transfer itself. Transfers that depend on income or spending change behaviour
and so cost efficiency. The size of that cost, and how it compares with the value of a different
distribution, is a question for the public economics of the Year 3 volume, not for these theorems.

\section{Exercises}\label{sec:b2:welfare-theorems:exercises}

\begin{exercise}[$\star$]\label{exo:b2:welfare-theorems:1}
Crusoe has $T = 16$ hours, produces $c = 4\sqrt L$ and has $u = \ln c + \ln(16 - L)$. Find the efficient
hours, the real wage that supports them, the firm's profit, and check that the household chooses the same
hours at that wage.
\end{exercise}

\begin{exercise}[$\star$]\label{exo:b2:welfare-theorems:2}
Read the opening chart. By how much did the price gap fall between the first and the last period, and by
how much did freight fall? What share of the fall does freight account for? What happened to the
residual?
\end{exercise}

\begin{exercise}[$\star$]\label{exo:b2:welfare-theorems:3}
A firm produces $y = 2\sqrt L$ and sells at $p = 3$; the wage is $w = 1$. Find its labour demand, output
and profit.
\end{exercise}

\begin{exercise}[$\star$]\label{exo:b2:welfare-theorems:4}
Describe price data from two connected markets that would refute \cref{prop:b2:welfare-theorems:one-price},
and say what else you would need to measure before concluding that it is refuted.
\end{exercise}

\begin{exercise}[$\star\star$]\label{exo:b2:welfare-theorems:5}
Crusoe's utility is $\tfrac23\ln c + \tfrac13\ln(16 - L)$, with $c = 4\sqrt L$. Find the efficient hours,
consumption and the supporting real wage. Show that in general, with $u = a\ln c + (1 - a)\ln(T - L)$ and
$c = A\sqrt L$, $L = aT/(2 - a)$.
\end{exercise}

\begin{exercise}[$\star\star$]\label{exo:b2:welfare-theorems:6}
In \cref{exo:b2:welfare-theorems:1}, productivity rises from $A = 4$ to $A = 5$. Predict, then compute, the
change in hours, consumption and the real wage. Why do hours not change?
\end{exercise}

\begin{exercise}[$\star\star$]\label{exo:b2:welfare-theorems:7}
Which assumption of the first welfare theorem fails in each case? (a) A factory's smoke soils the washing
of a laundry next door. (b) One company owns the only water source of a town. (c) Farmers cannot buy
insurance against a poor harvest. (d) A bridge costs a large sum to build and almost nothing to cross.
\end{exercise}

\begin{exercise}[$\star\star$]\label{exo:b2:welfare-theorems:8}
In the running economy of \cref{ch:b2:exchange-economies}, a planner wants $A$ to have 6 units of good 1.
Find the efficient \omterm{def:b2:exchange-economies:allocation}{allocation}, the supporting price and the transfer of good 2 that delivers it; check
$B$'s demand at that price.
\end{exercise}

\begin{exercise}[$\star\star$]\label{exo:b2:welfare-theorems:9}
Farmers in a region have no way to insure their harvests. Describe evidence that would show this missing
market makes the \omterm{def:b2:exchange-economies:allocation}{allocation} inefficient, rather than reflecting a risk that no one would insure at any
price.
\end{exercise}

\begin{exercise}[$\star\star\star$]\label{exo:b2:welfare-theorems:10}
Prove the first welfare theorem for a \omterm{def:b2:welfare-theorems:production-economy}{private-ownership economy} with production: if $(p, x^*, y^*)$ is a
\omterm{def:b2:exchange-economies:competitive-equilibrium}{competitive equilibrium} with $p \gg 0$ and preferences are \omterm{def:b2:welfare-theorems:lns}{locally non-satiated}, no \omterm{def:b2:exchange-economies:feasible}{feasible allocation}
$(x', y')$ Pareto-dominates $(x^*, y^*)$.
\end{exercise}

\begin{exercise}[$\star\star\star$]\label{exo:b2:welfare-theorems:11}
Two consumers have \omterm{def:b2:exchange-economies:endowment}{endowments} $(3, 3)$ each. $B$ has $u_B = x_{B1} + x_{B2}$; $A$ has $u_A = \min(x_{A1} +
x_{A2}, 4)$. Show that $p = (1, 1)$ with no trade is a \omterm{def:b2:exchange-economies:competitive-equilibrium}{competitive equilibrium}, that it is not
\omterm{def:b2:exchange-economies:pareto-efficient}{Pareto-efficient}, and say which step of the proof of \cref{thm:b2:welfare-theorems:first} fails.
\end{exercise}

\begin{exercise}[$\star\star\star$]\label{exo:b2:welfare-theorems:12}
Two Cobb--Douglas consumers spend shares $a$ and $b$ of their wealth on good 1; total \omterm{def:b2:exchange-economies:endowment}{endowments} are
$(\bar\omega_1, \bar\omega_2)$ and $A$ holds $\omega_A$. Derive the transfer of good 2 to $A$ that
delivers the efficient \omterm{def:b2:exchange-economies:allocation}{allocation} in which $A$ consumes $x_{A1}$ of good 1. Apply it with $a = \tfrac12$,
$b = \tfrac14$, $\bar\omega = (100, 100)$, $\omega_A = (80, 10)$ and $x_{A1} = 40$.
\end{exercise}

\section{Problem: Brenland's railway}\label{sec:b2:welfare-theorems:problem}

\begin{problem}[{Weekend problem --- a railway joins Ardel, a farming town of Brenland, an invented
economy, to the city of Brenmoor: what happens to the price of grain in each town, who gains, and what
transfer would deliver the allocation the council prefers}]
\label{pb:b2:welfare-theorems:1}
Ardel's farms grow grain; Ardel's households demand $100 - 10p_A$ tonnes a week at the price $p_A$ (in
dalens per tonne). Brenmoor grows none and demands $200 - 10p_B$. Grain moves by road from Ardel to
Brenmoor at a cost of $t = 6$ dalens a tonne, paid to competitive carriers. Round prices to two decimals.

\medskip
\noindent\textbf{Part I --- Two towns.}
\begin{enumerate}
  \item State the relation between $p_A$ and $p_B$ if grain is shipped, and if it is not.
  \item Farm output is $10p_A$ (Part II derives it). Write the market-clearing condition in each town,
        with $s$ the tonnes shipped.
  \item Solve for $p_A$, $p_B$ and $s$.
  \item Without the road, what would the price be in Ardel? What would Brenmoor consume?
  \item Find the freight bill.
  \item Given the cost of transport, is this \omterm{def:b2:exchange-economies:allocation}{allocation} efficient? Which theorem says so, and under
        which assumptions?
\end{enumerate}

\medskip
\noindent\textbf{Part II --- Production in Ardel.} Farms produce $q = \sqrt{20L}$ tonnes with $L$ hours
of labour, paid one dalen an hour.
\begin{enumerate}[resume]
  \item Derive the supply $q = 10p_A$ from profit maximisation.
  \item Find output, hours and profit at the price of question 3.
  \item Who receives the profit, and where does it appear in the households' budgets?
  \item Compute consumer surplus in each town and the farms' profit, and their total.
\end{enumerate}

\medskip
\noindent\textbf{Part III --- The railway.} The railway cuts the transport cost to $t = 2$.
\begin{enumerate}[resume]
  \item Find the new prices.
  \item Find the tonnes shipped, Ardel's output and Ardel's own consumption.
  \item Check the gap between the two prices.
  \item Compute the new surpluses and profit, and the gain in total surplus.
  \item Who loses? By how much?
  \item Could the losers be compensated with the gain? What kind of transfer would do it without
        changing anyone's choices?
  \item In nineteenth-century data the gap between two markets exceeded the cost of shipping. What would a
        positive residual mean in this model?
  \item What evidence after the railway opens would show that the two markets had become one?
\end{enumerate}

\medskip
\noindent\textbf{Part IV --- A transfer.} Once a week, Ardel and Brenmoor also trade grain against cloth
(the \omterm{def:b2:exchange-economies:numeraire}{numeraire}). Ardel holds 80 tonnes of grain and 10 bolts of cloth, Brenmoor 20 tonnes and 90 bolts;
Ardel spends half its wealth on grain and Brenmoor a quarter (Cobb--Douglas).
\begin{enumerate}[resume]
  \item Find the \omterm{def:b2:exchange-economies:competitive-equilibrium}{competitive equilibrium} price of grain in cloth and the \omterm{def:b2:exchange-economies:allocation}{allocation}.
  \item The council wants Brenmoor to consume 60 tonnes of grain, at an efficient \omterm{def:b2:exchange-economies:allocation}{allocation}. Find
        Ardel's cloth at that \omterm{def:b2:exchange-economies:allocation}{allocation}.
  \item Find the price that supports it.
  \item Find the transfer of cloth between the towns that delivers it.
  \item Check by solving the equilibrium after the transfer.
  \item State the result: the prices in the two towns after the railway, and the transfer that delivers
        the council's \omterm{def:b2:exchange-economies:allocation}{allocation}.
\end{enumerate}
\end{problem}
